\section{Introduction}

Dynamic discrete choice models in economics and inverse reinforcement learning in computer science solve the same fundamental inference problem. An agent makes a sequence of decisions over time, and the researcher observes these decisions and wants to recover the reward function that rationalizes them. \citet{rust1987} formalized this problem as nested fixed-point maximum likelihood estimation, and the structural econometrics literature that followed developed a rich family of estimators including conditional choice probability methods and mathematical programming approaches \citep{aguirregabiriaMira2010}. Independently, \citet{ziebart2010} formalized the same problem as maximum causal entropy inverse reinforcement learning, and the machine learning literature developed its own family of estimators including adversarial methods and soft Q-learning. These two traditions share the same Bellman equation, the same discount factor, and the same logit choice kernel, yet they developed with different notation, different software ecosystems, and no unified pipeline for statistical inference or counterfactual analysis.

Existing software does not bridge these traditions. The \texttt{imitation} library \citep{gleave2022} implements several IRL algorithms but provides no standard errors, no identification diagnostics, and no counterfactual simulation. Stata packages for dynamic discrete choice cover NFXP and a handful of related estimators but cannot accommodate neural reward functions or adversarial training. There is no library where a practitioner can specify a DDC model once, swap between structural and IRL estimators, and obtain valid inference from all of them through a single API. This gap forces applied researchers to maintain separate codebases for each estimator, making systematic comparison across methods unnecessarily difficult.

This paper introduces EconIRL, a Python library that provides twelve estimators for dynamic discrete choice models through a common \texttt{Estimator} protocol. Six estimators come from structural econometrics (NFXP, CCP, MPEC, NNES, SEES, TD-CCP) and six from inverse reinforcement learning (MCE-IRL, Deep MaxEnt, AIRL, GLADIUS, IQ-Learn, Behavioral Cloning). Every estimator returns an \texttt{EstimationSummary} object that provides standard errors by four methods (analytical, bootstrap, jackknife, and sandwich), identification diagnostics, and counterfactual simulation. The library is built on JAX \citep{jax2018} for automatic differentiation through fixed-point solvers via implicit differentiation, with Equinox \citep{kidger2021} for neural network components. We also present a new theoretical result connecting the IQ-Learn chi-squared divergence objective to a penalized conditional log-likelihood in the tabular DDC setting \citep{rawat2026}, which provides a formal bridge between occupancy-matching IRL and structural maximum likelihood estimation. Three simulation studies validate convergence, scaling, and nonlinear reward recovery across estimator families. The library is MIT licensed and pip-installable.

The remainder of the paper is organized as follows. Section~\ref{sec:framework} sets up notation and describes the dynamic discrete choice model that serves as the common foundation for all estimators. Section~\ref{sec:estimators} describes each estimator with pseudocode and references to the underlying theory. Section~\ref{sec:iq_learn} presents the IQ-Learn equivalence result. Section~\ref{sec:post_estimation} covers the post-estimation infrastructure for inference, diagnostics, and counterfactual analysis. Section~\ref{sec:experiments} presents three simulation case studies and real-data applications. Section~\ref{sec:conclusion} concludes.
